% GENERATED FILE. Edit canonical lessons, then run npm run generate:books.
% canonical-source-hash: fnv1a64:438319202cc9fb22
% canonical-lessons: JA-C23-kusuri, JA-C23-nemui, JA-C23-kanashii, JA-C23-tanoshii, JA-C23-kowai

\chapter{Medicine, Sleepy, Sad, Fun, Scary}
\label{ch:ja-medicine-sleepy-sad-fun-scary}
\hlchaptermodality{23}

\begin{chapteropening}
\textbf{By the end of this chapter:} \emph{I can say medicine, sleepy, sad, fun, enjoyable and scary.}

Complete the last lesson of chapter 23: i can say medicine, sleepy, sad, fun, enjoyable and scary.
\end{chapteropening}

\section[kusuri]{\ja{くすり} (kusuri) — medicine}

\label{lesson:JA-C23-kusuri}



\begin{quote}
Before the new one: say the Japanese for a fever, then the Japanese for painful.
\end{quote}

\subsection*{You'll want to know: \ja{くすり}}
\textbf{\ja{くすり}} — \emph{kusuri} — \textquotedblleft{}medicine\textquotedblright{}.

\begin{grammarlens}[title={What you've built}]
One more word for this chapter.
\end{grammarlens}

\subsection*{Your turn}
\begin{itemize}
\raggedright
  \item \emph{Say it:} \emph{kusuri}
  \item \emph{Say it:} \emph{kusuri}, once more
  \item \emph{Say it:} say \emph{kusuri}
  \item \emph{Recall:} say the Japanese for a fever, then the Japanese for painful, then say \emph{kusuri} again
\end{itemize}

\begin{tcolorbox}[breakable,colback=teal!4,colframe=teal!35!black,title={Before you move on}]
What does \ja{くすり} mean? (\textquotedblleft{}Medicine\textquotedblright{}.) Say it once more.
\end{tcolorbox}
\section[nemui]{\ja{ねむい} (nemui) — sleepy}

\label{lesson:JA-C23-nemui}



\begin{quote}
Before the new one: say the Japanese for painful, then the Japanese for medicine.
\end{quote}

\subsection*{You'll want to know: \ja{ねむい}}
\textbf{\ja{ねむい}} — \emph{nemui} — \textquotedblleft{}sleepy\textquotedblright{}.

\begin{grammarlens}[title={What you've built}]
One more word for this chapter.
\end{grammarlens}

\subsection*{Your turn}
\begin{itemize}
\raggedright
  \item \emph{Say it:} \emph{nemui}
  \item \emph{Say it:} \emph{nemui}, once more
  \item \emph{Say it:} say \emph{nemui}
  \item \emph{Recall:} say the Japanese for painful, then the Japanese for medicine, then say \emph{nemui} again
\end{itemize}

\begin{tcolorbox}[breakable,colback=teal!4,colframe=teal!35!black,title={Before you move on}]
What does \ja{ねむい} mean? (\textquotedblleft{}Sleepy\textquotedblright{}.) Say it once more.
\end{tcolorbox}
\section[kanashii]{\ja{かなしい} (kanashii) — sad}

\label{lesson:JA-C23-kanashii}



\begin{quote}
Before the new one: say the Japanese for medicine, then the Japanese for sleepy.
\end{quote}

\subsection*{You'll want to know: \ja{かなしい}}
\textbf{\ja{かなしい}} — \emph{kanashii} — \textquotedblleft{}sad\textquotedblright{}.

\begin{grammarlens}[title={What you've built}]
One more word for this chapter.
\end{grammarlens}

\subsection*{Your turn}
\begin{itemize}
\raggedright
  \item \emph{Say it:} \emph{kanashii}
  \item \emph{Say it:} \emph{kanashii}, once more
  \item \emph{Say it:} say \emph{kanashii}
  \item \emph{Recall:} say the Japanese for medicine, then the Japanese for sleepy, then say \emph{kanashii} again
\end{itemize}

\begin{tcolorbox}[breakable,colback=teal!4,colframe=teal!35!black,title={Before you move on}]
What does \ja{かなしい} mean? (\textquotedblleft{}Sad\textquotedblright{}.) Say it once more.
\end{tcolorbox}
\section[tanoshii]{\ja{たのしい} (tanoshii) — fun, enjoyable}

\label{lesson:JA-C23-tanoshii}



\begin{quote}
Before the new one: say the Japanese for sleepy, then the Japanese for sad.
\end{quote}

\subsection*{You'll want to know: \ja{たのしい}}
\textbf{\ja{たのしい}} — \emph{tanoshii} — \textquotedblleft{}fun, enjoyable\textquotedblright{}.

\begin{grammarlens}[title={What you've built}]
One more word for this chapter.
\end{grammarlens}

\subsection*{Your turn}
\begin{itemize}
\raggedright
  \item \emph{Say it:} \emph{tanoshii}
  \item \emph{Say it:} \emph{tanoshii}, once more
  \item \emph{Say it:} say \emph{tanoshii}
  \item \emph{Recall:} say the Japanese for sleepy, then the Japanese for sad, then say \emph{tanoshii} again
\end{itemize}

\begin{tcolorbox}[breakable,colback=teal!4,colframe=teal!35!black,title={Before you move on}]
What does \ja{たのしい} mean? (\textquotedblleft{}Fun, enjoyable\textquotedblright{}.) Say it once more.
\end{tcolorbox}
\section[kowai]{\ja{こわい} (kowai) — scary}

\label{lesson:JA-C23-kowai}



\begin{quote}
Before the new one: say the Japanese for sad, then the Japanese for fun.
\end{quote}

\subsection*{You'll want to know: \ja{こわい}}
\textbf{\ja{こわい}} — \emph{kowai} — \textquotedblleft{}scary\textquotedblright{}.

\begin{grammarlens}[title={What you've built}]
One more word for this chapter.
\end{grammarlens}

\subsection*{Your turn}
\begin{itemize}
\raggedright
  \item \emph{Say it:} \emph{kowai}
  \item \emph{Say it:} \emph{kowai}, once more
  \item \emph{Say it:} say \emph{kowai}
  \item \emph{Recall:} say the Japanese for sad, then the Japanese for fun, then say \emph{kowai} again
\end{itemize}

\begin{tcolorbox}[breakable,colback=teal!4,colframe=teal!35!black,title={Before you move on}]
What does \ja{こわい} mean? (\textquotedblleft{}Scary\textquotedblright{}.) Say it once more.
\end{tcolorbox}
